\section{Hamilton 力学}\label{sec:hamiltonian-mechanics}

下面将看到, 在一定条件下, 可以把 Euler–Lagrange 方程作为余切丛上具有特殊几何性质的流来研究.

设 $M$ 为 $n$ 维流形. 集合
\[
  TM\oplus T^*M:=\bigcup_{q\in M}T_qM\times T_q^*M
\]
具有一个显然的可微结构: 若 $(x^1,\ldots,x^n)$ 为 $M$ 上的局部坐标, 则 $TM\oplus T^*M$ 上标记 $(v,\omega)\in T_qM\times T_q^*M$ 的局部坐标为 $(x^1,\ldots,x^n,v^1,\ldots,v^n,p_1,\ldots,p_n)$, 其中
\[
  v=\sum_{i=1}^n v^i\frac{\partial}{\partial x^i},
  \qquad \omega=\sum_{i=1}^n p_i\,\mathrm{d}x^i,
\]
且 $q$ 是坐标为 $(x^1,\ldots,x^n)$ 的点. 对于这个可微结构, 由 $\pi_1(v,\omega)=v$ 和 $\pi_2(v,\omega)=\omega$ 给出的映射 $\pi_1:TM\oplus T^*M\to TM$ 和 $\pi_2:TM\oplus T^*M\to T^*M$ 都是浸没.

\begin{RiemannDefinition}\label{def:extended-hamiltonian-function}
对应于 Lagrange 函数 $L:TM\to\mathbb{R}$ 的\textbf{扩展 Hamilton 函数}是映射 $\widetilde H:TM\oplus T^*M\to\mathbb{R}$, 由
\[
  \widetilde H(v,\omega):=\omega(v)-L(v)
\]
给出.
\end{RiemannDefinition}

在局部坐标下有
\begin{align*}
  \widetilde H(x^1,\ldots,x^n,v^1,\ldots,v^n,p_1,\ldots,p_n)
  &=\sum_{i=1}^n p_iv^i\\
  &\quad-L(x^1,\ldots,x^n,v^1,\ldots,v^n),
\end{align*}
从而
\[
  \mathrm{d}\widetilde H
  =\sum_{i=1}^n\left(p_i-\frac{\partial L}{\partial v^i}\right)\mathrm{d}v^i
   +\sum_{i=1}^n v^i\,\mathrm{d}p_i
   -\sum_{i=1}^n\frac{\partial L}{\partial x^i}\,\mathrm{d}x^i.
\]
因此, 将 $\widetilde H$ 限制在任意形如 $T_qM\times\{\omega\}$ 的子流形上 (固定 $q\in M$ 和 $\omega\in T_q^*M$), 所得函数的任一临界点必满足
\[
  p_i=\frac{\partial L}{\partial v^i}(x^1,\ldots,x^n,v^1,\ldots,v^n)
  \qquad(i=1,\ldots,n).
\]
由此可知, 所有这些临界点自然组成一个 $2n$ 维子流形 $S\subset TM\oplus T^*M$, 使 $\pi_1|_S:S\to TM$ 为微分同胚. 若 $\pi_2|_S:S\to T^*M$ 也是微分同胚, 则称该 Lagrange 函数\textbf{超正则}. 在此情形, $\pi_2|_S\circ(\pi_1|_S)^{-1}:TM\to T^*M$ 是一个保纤维的微分同胚, 称为 \textbf{Legendre 变换}.

给定超正则的 Lagrange 函数, 可以用映射 $\pi_1|_S$ 和 $\pi_2|_S$ 作等同 $TM\cong S\cong T^*M$. Hamilton 函数 $H:TM\to\mathbb{R}$ 与扩展 Hamilton 函数由 $H=\widetilde H\circ(\pi_1|_S)^{-1}$ 相联系, 因此在这些等同下可直接写成 $H=\widetilde H|_S$. 所以
\[
  \mathrm{d}H=\sum_{i=1}^n v^i\,\mathrm{d}p_i
  -\sum_{i=1}^n\frac{\partial L}{\partial x^i}\,\mathrm{d}x^i
\]
(这里必须把 $x^1,\ldots,x^n,v^1,\ldots,v^n,p_1,\ldots,p_n$ 看作 $S$ 上的局部函数, 使 $(x^1,\ldots,x^n,v^1,\ldots,v^n)$ 和 $(x^1,\ldots,x^n,p_1,\ldots,p_n)$ \textbf{都}是局部坐标). 另一方面, 把 $H$ 看作余切丛上的函数, 得到
\[
  \mathrm{d}H=\sum_{i=1}^n\frac{\partial H}{\partial x^i}\,\mathrm{d}x^i
  +\sum_{i=1}^n\frac{\partial H}{\partial p_i}\,\mathrm{d}p_i.
\]
因此必有
\[
  \begin{cases}
     \dfrac{\partial H}{\partial x^i}=-\dfrac{\partial L}{\partial x^i}\\[.6ex]
     \dfrac{\partial H}{\partial p_i}=v^i
  \end{cases}
  \qquad(i=1,\ldots,n),
\]
其中 Hamilton 函数的偏导数须相对于局部坐标 $(x^1,\ldots,x^n,p_1,\ldots,p_n)$ 计算, 而 Lagrange 函数的偏导数须相对于局部坐标 $(x^1,\ldots,x^n,v^1,\ldots,v^n)$ 计算.

\begin{RiemannProposition}\label{prop:legendre-conjugates-euler-lagrange-and-hamilton}
超正则 Lagrange 函数 $L:TM\to\mathbb{R}$ 的 Euler–Lagrange 方程在 $TM$ 上定义一个流. 通过 Legendre 变换, 这个流被映为 $T^*M$ 上由 \textbf{Hamilton 方程}
\[
  \begin{cases}
    \dot x^i=\dfrac{\partial H}{\partial p_i}\\[.6ex]
    \dot p_i=-\dfrac{\partial H}{\partial x^i}
  \end{cases}
  \qquad(i=1,\ldots,n)
\]
所定义的流.
\end{RiemannProposition}

\begin{RiemannProof}
Euler–Lagrange 方程可以写成 $TM$ 上的下列一阶常微分方程组:
\[
  \begin{cases}
    \dot x^i=v^i\\[.6ex]
    \dfrac{\mathrm{d}}{\mathrm{d}t}\left(\dfrac{\partial L}{\partial v^i}\right)
    =\dfrac{\partial L}{\partial x^i}
  \end{cases}
  \qquad(i=1,\ldots,n).
\]
由于在 $S$ 上有
\[
  p_i=\frac{\partial L}{\partial v^i},\qquad
  v^i=\frac{\partial H}{\partial p_i},\qquad
  \frac{\partial L}{\partial x^i}=-\frac{\partial H}{\partial x^i},
\]
可见在局部坐标 $(x^1,\ldots,x^n,p_1,\ldots,p_n)$ 下, 这个方程组化为 Hamilton 方程. Hamilton 方程显然在 $T^*M$ 上定义一个流, 因此 Euler–Lagrange 方程必在 $TM$ 上定义一个流.
\end{RiemannProof}

\begin{RiemannExample}\label{ex:legendre-transform-for-conservative-system}
保守力学系统 $(M,\langle\cdot,\cdot\rangle,-\mathrm{d}U)$ 的 Lagrange 函数在局部坐标下写为
\[
  L(x^1,\ldots,x^n,v^1,\ldots,v^n)
  =\frac{1}{2}\sum_{i,j=1}^n g_{ij}(x^1,\ldots,x^n)v^iv^j
  -U(x^1,\ldots,x^n).
\]
在这些坐标下, Legendre 变换由
\[
  p_i=\frac{\partial L}{\partial v^i}=\sum_{j=1}^n g_{ij}v^j
  \qquad(i=1,\ldots,n)
\]
给出, 它确为保纤维的微分同胚, 其逆由
\[
  v^i=\sum_{j=1}^n g^{ij}p_j\qquad(i=1,\ldots,n)
\]
给出. 作为切丛上的函数, Hamilton 函数为 (参见 Example~\ref{ex:hamiltonian-as-mechanical-energy})
\[
  H=\frac{1}{2}\sum_{i,j=1}^n g_{ij}v^iv^j+U.
\]
利用 Legendre 变换, 可将 Hamilton 函数视为余切丛上的下列函数:
\[
  H=\frac{1}{2}\sum_{i,j,k,l=1}^n g_{ij}g^{ik}p_kg^{jl}p_l+U
    =\frac{1}{2}\sum_{k,l=1}^n g^{kl}p_kp_l+U.
\]
因此保守力学系统的 Hamilton 方程为
\[
  \begin{cases}
    \displaystyle\dot x^i=\sum_{j=1}^n g^{ij}p_j\\[.6ex]
    \displaystyle\dot p_i=-\frac12\sum_{k,l=1}^n
       \frac{\partial g^{kl}}{\partial x^i}p_kp_l-\frac{\partial U}{\partial x^i}
  \end{cases}
  \qquad(i=1,\ldots,n).
\]
\end{RiemannExample}

Hamilton 方程定义的流具有一些显著的几何性质. 引入下面的定义后, 这些性质会更容易理解.

\begin{RiemannDefinition}\label{def:canonical-symplectic-potential-and-form}
\textbf{典范辛势} (或 \textbf{Liouville 形式}) 是 1-形式 $\theta\in\Omega^1(T^*M)$, 对所有 $\alpha\in T^*M$ 和 $v\in T_\alpha(T^*M)$ 由
\[
  \theta_\alpha(v):=\alpha\bigl((\mathrm{d}\pi)_\alpha(v)\bigr)
\]
给出, 其中 $\pi:T^*M\to M$ 为自然投影. $T^*M$ 上的\textbf{典范辛形式}是由 $\omega=\mathrm{d}\theta$ 给出的 2-形式 $\omega\in\Omega^2(T^*M)$.
\end{RiemannDefinition}

在局部坐标下,
\[
  \pi(x^1,\ldots,x^n,p_1,\ldots,p_n)=(x^1,\ldots,x^n)
\]
且
\[
  v=\sum_{i=1}^n\mathrm{d}x^i(v)\frac{\partial}{\partial x^i}
  +\sum_{i=1}^n\mathrm{d}p_i(v)\frac{\partial}{\partial p_i}.
\]
因此
\[
  (\mathrm{d}\pi)_\alpha(v)
  =\sum_{i=1}^n\mathrm{d}x^i(v)\frac{\partial}{\partial x^i},
\]
从而
\[
  \theta_\alpha(v)=\alpha\bigl((\mathrm{d}\pi)_\alpha(v)\bigr)
  =\sum_{i=1}^n p_i\,\mathrm{d}x^i\left(
     \sum_{j=1}^n\mathrm{d}x^j(v)\frac{\partial}{\partial x^j}\right)
  =\sum_{i=1}^n p_i\,\mathrm{d}x^i(v).
\]
由此得到
\[
  \theta=\sum_{i=1}^n p_i\,\mathrm{d}x^i,
\]
进而
\[
  \omega=\sum_{i=1}^n\mathrm{d}p_i\wedge\mathrm{d}x^i.
\]

\begin{RiemannProposition}\label{prop:closed-nondegenerate-canonical-symplectic-form}
典范辛形式 $\omega$ 是闭的 ($\mathrm{d}\omega=0$) 且非退化. 此外, $\omega^n=\omega\wedge\cdots\wedge\omega$ 是体积形式 (特别地, 即使 $M$ 本身不可定向, $T^*M$ 也总是可定向的).
\end{RiemannProposition}

\begin{RiemannProof}
Exercise~\ref{ex:hamiltonian-mechanics}(\ref{ex:verify-canonical-symplectic-form}).
\end{RiemannProof}

回顾第~\ref{chap:differential-forms}~章的 Exercise~\ref{ex:tensors}(\ref{ex:contraction-of-alternating-tensors}): 若 $v\in T_pM$, 则 $\iota(v)\omega\in T_p^*M$ 是对所有 $w\in T_pM$ 由
\[
  (\iota(v)\omega)(w)=\omega(v,w)
\]
给出的余向量. 因此 Proposition~\ref{prop:closed-nondegenerate-canonical-symplectic-form} 的第一个论断等价于说, 对所有 $p\in M$, 映射 $T_pM\ni v\mapsto\iota(v)\omega\in T_p^*M$ 都是线性同构.

下面的结果给出了理解 Hamilton 方程几何意义的关键.

\begin{RiemannProposition}\label{prop:hamilton-equations-as-symplectic-duality}
Hamilton 方程是满足
\[
  \iota(X_H)\omega=-\mathrm{d}H
\]
的向量场 $X_H$ 的流的方程.
\end{RiemannProposition}

\begin{RiemannProof}
Hamilton 方程给出向量场
\[
  X_H=\sum_{i=1}^n\left(
    \frac{\partial H}{\partial p_i}\frac{\partial}{\partial x^i}
   -\frac{\partial H}{\partial x^i}\frac{\partial}{\partial p_i}\right)
\]
的流. 因此
\begin{align*}
  \iota(X_H)\omega
  &=\iota(X_H)\sum_{i=1}^n(\mathrm{d}p_i\otimes\mathrm{d}x^i
                        -\mathrm{d}x^i\otimes\mathrm{d}p_i)\\
  &=\sum_{i=1}^n\left(-\frac{\partial H}{\partial x^i}\,\mathrm{d}x^i
                    -\frac{\partial H}{\partial p_i}\,\mathrm{d}p_i\right)
    =-\mathrm{d}H.
\end{align*}
\end{RiemannProof}

\begin{RiemannRemark}\label{rem:symplectic-gradient-of-hamiltonian}
注意, 由于 $\omega$ 非退化, $H$ 完全确定 $X_H$. 类比 Riemann 情形, $-X_H$ 有时称为 $H$ 的\textbf{辛梯度}. 向量场 $X_H$ 通常称为由 $H$ 确定的 \textbf{Hamilton 向量场}.
\end{RiemannRemark}

\begin{RiemannDefinition}\label{def:hamiltonian-flow}
由 $F\in C^\infty(T^*M)$ 生成的 \textbf{Hamilton 流}是满足
\[
  \iota(X_F)\omega=-\mathrm{d}F
\]
的唯一向量场 $X_F\in\mathfrak{X}(T^*M)$ 的流.
\end{RiemannDefinition}

因此, 超正则 Lagrange 函数在 $T^*M$ 上确定的流是 Hamilton 流的一种特殊情形 (此时生成函数为 Hamilton 函数). 下面讨论一般 Hamilton 流的几何性质, 先从 Hamilton 形式的能量守恒讲起.

\begin{RiemannProposition}\label{prop:hamiltonian-flow-preserves-generator}
Hamilton 流保持其生成函数.
\end{RiemannProposition}

\begin{RiemannProof}
由于 $\omega$ 是交替的, 有
\[
  X_F\cdot F=\mathrm{d}F(X_F)=(-\iota(X_F)\omega)(X_F)=-\omega(X_F,X_F)=0.
\]
\end{RiemannProof}

\begin{RiemannProposition}\label{prop:hamiltonian-flow-preserves-symplectic-form}
Hamilton 流保持典范辛形式: 若 $\psi_t:T^*M\to T^*M$ 是 Hamilton 流, 则 $\psi_t^*\omega=\omega$.
\end{RiemannProposition}

\begin{RiemannProof}
设 $F\in C^\infty(T^*M)$ 为其 Hamilton 流是 $\psi_t$ 的函数. 回顾第~\ref{chap:differential-forms}~章的 Exercise~\ref{ex:differential-forms}(\ref{ex:lie-derivative-of-differential-form}): $\omega$ 沿 $X_F\in\mathfrak{X}(T^*M)$ 的 Lie 导数
\[
  L_{X_F}\omega=\left.\frac{\mathrm{d}}{\mathrm{d}t}\right|_{t=0}\psi_t^*\omega
\]
可由 Cartan 公式计算:
\[
  L_{X_F}\omega=\iota(X_F)\mathrm{d}\omega
  +\mathrm{d}(\iota(X_F)\omega)=\mathrm{d}(-\mathrm{d}F)=0.
\]
因此
\begin{align*}
  \frac{\mathrm{d}}{\mathrm{d}t}\psi_t^*\omega
  &=\left.\frac{\mathrm{d}}{\mathrm{d}s}\right|_{s=0}(\psi_{t+s})^*\omega
   =\left.\frac{\mathrm{d}}{\mathrm{d}s}\right|_{s=0}(\psi_s\circ\psi_t)^*\omega
   =\left.\frac{\mathrm{d}}{\mathrm{d}s}\right|_{s=0}\psi_t^*(\psi_s^*\omega)\\
  &=\psi_t^*\left.\frac{\mathrm{d}}{\mathrm{d}s}\right|_{s=0}\psi_s^*\omega
   =\psi_t^*L_{X_F}\omega=0.
\end{align*}
由此得出
\[
  \psi_t^*\omega=(\psi_0)^*\omega=\omega.
\]
\end{RiemannProof}

\begin{RiemannTheorem}[Liouville]\label{thm:liouville-preservation-of-symplectic-volume}
Hamilton 流保持关于辛体积形式的积分: 若 $\psi_t:T^*M\to T^*M$ 是 Hamilton 流, $F\in C^\infty(T^*M)$ 是紧支集函数, 则
\[
  \int_{T^*M}F\circ\psi_t=\int_{T^*M}F.
\]
\end{RiemannTheorem}

\begin{RiemannProof}
这是 $\psi_t$ 保持辛体积形式的直接推论, 因为
\[
  \psi_t^*(\omega^n)=(\psi_t^*\omega)^n=\omega^n.
\]
因此
\begin{align*}
  \int_{T^*M}F\circ\psi_t
  &=\int_{T^*M}(F\circ\psi_t)\omega^n
   =\int_{T^*M}(F\circ\psi_t)\psi_t^*(\omega^n)\\
  &=\int_{T^*M}\psi_t^*(F\omega^n)=\int_{T^*M}F\omega^n=\int_{T^*M}F
\end{align*}
[参见第~\ref{chap:differential-forms}~章的 Exercise~\ref{ex:integration-on-manifolds}(\ref{ex:orientation-preserving-pullback-integral})].
\end{RiemannProof}

\begin{RiemannTheorem}[Poincaré 回归]\label{thm:poincare-recurrence-for-hamiltonian-flow}
设 $\psi_t:T^*M\to T^*M$ 为 Hamilton 流, $K\subset T^*M$ 为在 $\psi_t$ 下不变的紧集. 则对于每个开集 $U\subset K$ 和每个 $T>0$, 都存在 $\alpha\in U$ 和 $t\geq T$, 使得 $\psi_t(\alpha)\in U$.
\end{RiemannTheorem}

\begin{RiemannProof}
设 $F\in C^\infty(T^*M)$ 是取值于 $[0,1]$ 的紧支集光滑函数, 使所有 $\alpha\in K$ 都满足 $F(\alpha)=1$ (例如, 利用单位分解容易构造这样的函数). 设 $G\in C^\infty(T^*M)$ 是取值于 $[0,1]$, 且紧支集包含在 $U$ 中的光滑函数, 满足
\[
  \int_{T^*M}G>0.
\]
考虑开集 $U_n:=\psi_{nT}(U)$. 若这些集合两两不交, 则可对每个 $N\in\mathbb{N}$ 定义函数 $\widetilde G_N\in C^\infty(M)$:
\[
  \widetilde G_N(\alpha)=
  \begin{cases}
    (G\circ\psi_{-nT})(\alpha) & \text{若 $\alpha\in U_n$ 且 $n\leq N$}\\
    0 & \text{其他情形}
  \end{cases}.
\]
这些函数的紧支集包含在 $K$ 中 (因为 $K$ 在 $\psi_t$ 下不变), 取值于 $[0,1]$, 因而满足 $\widetilde G_N\leq F$. 因此, 对所有 $N\in\mathbb{N}$, 都有
\[
  \int_{T^*M}F\geq\int_{T^*M}\widetilde G_N
   =\sum_{n=1}^N\int_{T^*M}G\circ\psi_{-nT}=N\int_{T^*M}G,
\]
这是不可能的. 由此可知必存在 $m,n\in\mathbb{N}$ (不妨设 $n>m$), 使得
\[
  U_m\cap U_n\ne\varnothing
  \iff\psi_{mT}(U)\cap\psi_{nT}(U)\ne\varnothing
  \iff U\cap\psi_{(n-m)T}(U)\ne\varnothing.
\]
取 $t=(n-m)T$ 和 $\alpha\in\psi_{-t}(U\cap\psi_t(U))=\psi_{-t}(U)\cap U$, 即得所需结论.
\end{RiemannProof}

可利用余切丛的辛结构在 $T^*M$ 上的可微函数的集合上定义一种新的二元运算.

\begin{RiemannDefinition}\label{def:poisson-bracket}
两个可微函数 $F,G\in C^\infty(T^*M)$ 的 \textbf{Poisson 括号}为 $\{F,G\}:=X_F\cdot G$.
\end{RiemannDefinition}

\begin{RiemannProposition}\label{prop:poisson-bracket-lie-algebra}
$(C^\infty(T^*M),\{\cdot,\cdot\})$ 是一个 Lie 代数, 将函数 $F\in C^\infty(T^*M)$ 映到其 Hamilton 向量场 $X_F\in\mathfrak{X}(T^*M)$ 的映射是 Lie 代数同态, 即对任意 $F,G,H\in C^\infty(T^*M)$ 和 $\alpha,\beta\in\mathbb{R}$, 有:
\begin{enumerate}[label=(\roman*),leftmargin=*]
\item $\{F,G\}=-\{G,F\}$;
\item $\{\alpha F+\beta G,H\}=\alpha\{F,H\}+\beta\{G,H\}$;
\item $\{F,\{G,H\}\}+\{G,\{H,F\}\}+\{H,\{F,G\}\}=0$;
\item $X_{\{F,G\}}=[X_F,X_G]$.
\end{enumerate}
\end{RiemannProposition}

\begin{RiemannProof}
有
\begin{align*}
  \{F,G\}&=X_F\cdot G=\mathrm{d}G(X_F)=(-\iota(X_G)\omega)(X_F)\\
  &=-\omega(X_G,X_F)=\omega(X_F,X_G),
\end{align*}
这证明了 Poisson 括号的反对称性和双线性. 另一方面,
\begin{align*}
  \iota(X_{\{F,G\}})\omega
   &=-\mathrm{d}\{F,G\}
    =-\mathrm{d}(X_F\cdot G)
    =-\mathrm{d}(\mathrm{d}G(X_F))
    =-\mathrm{d}(\iota(X_F)\mathrm{d}G)\\
   &=-L_{X_F}\mathrm{d}G
    =L_{X_F}(\iota(X_G)\omega)
    =\iota(L_{X_F}X_G)\omega+\iota(X_G)L_{X_F}\omega\\
   &=\iota([X_F,X_G])\omega
\end{align*}
[参见第~\ref{chap:differential-forms}~章的 Exercise~\ref{ex:differential-forms}(\ref{ex:lie-derivative-of-differential-form})]. 由于 $\omega$ 非退化, 得到
\[
  X_{\{F,G\}}=[X_F,X_G].
\]
最后,
\begin{align*}
  &\{F,\{G,H\}\}+\{G,\{H,F\}\}+\{H,\{F,G\}\}\\
  &\quad=\{F,X_G\cdot H\}-\{G,X_F\cdot H\}-X_{\{F,G\}}\cdot H\\
  &\quad=X_F\cdot(X_G\cdot H)-X_G\cdot(X_F\cdot H)-[X_F,X_G]\cdot H=0.
\end{align*}
\end{RiemannProof}

\begin{RiemannExercise}\label{ex:hamiltonian-mechanics}
\begin{enumerate}[label=(\arabic*),ref=\arabic*,leftmargin=*]
\item\label{ex:verify-canonical-symplectic-form}
证明 Proposition~\ref{prop:closed-nondegenerate-canonical-symplectic-form}.
\item\label{ex:conservative-system-with-magnetic-term}
设 $(M,\langle\cdot,\cdot\rangle)$ 为 Riemann 流形, $\alpha\in\Omega^1(M)$ 为 1-形式, $U\in C^\infty(M)$ 为可微函数.
\begin{enumerate}[label=(\alph*),leftmargin=*]
\item 证明由下式对 $v\in T_pM$ 所定义的 Lagrange 函数 $L:TM\to\mathbb{R}$
\[
  L(v)=\frac{1}{2}\langle v,v\rangle+\iota(v)\alpha_p-U(p)
\]
的 Euler–Lagrange 方程给出力学系统 $(M,\langle\cdot,\cdot\rangle,\mathcal{F})$ 的运动, 其中对 $v\in T_pM$ 有
\[
  \mathcal{F}(v)=-(\mathrm{d}U)_p-\iota(v)(\mathrm{d}\alpha)_p.
\]
\item 证明机械能 $E=K+U$ 沿 $(M,\langle\cdot,\cdot\rangle,\mathcal{F})$ 的运动守恒 (因此称其为\textbf{带磁项的保守力学系统}).
\item 证明 $L$ 超正则, 并求 Legendre 变换.
\item 求 Hamilton 函数 $H:T^*M\to\mathbb{R}$, 并写出 Hamilton 方程.
\end{enumerate}
\item\label{ex:free-relativistic-particle-legendre-transform}
设正数 $c>0$ 表示光速, 并考虑开集 $U:=\{v\in T\mathbb{R}^n\mid\|v\|<c\}$, 其中 $\|\cdot\|$ 为 Euclidean 范数. 质量为 $m>0$ 的自由相对论粒子的运动由 Lagrange 函数 $L:U\to\mathbb{R}$ 确定, 其中
\[
  L(v):=-mc^2\sqrt{1-\frac{\|v\|^2}{c^2}}.
\]
\begin{enumerate}[label=(\alph*),leftmargin=*]
\item 证明 $L$ 超正则, 并求 Legendre 变换.
\item 求 Hamilton 函数 $H:T^*\mathbb{R}^n\to\mathbb{R}$, 并写出 Hamilton 方程.
\end{enumerate}
\item\label{ex:density-of-poincare-recurrent-points}
证明在 Poincaré 回归定理中, 对某个 $t\geq T$ 满足 $\psi_t(\alpha)\in U$ 的点 $\alpha\in U$ 构成的集合在 $U$ 中稠密. (\textbf{注}: 可以证明这是一个满测集.)
\item\label{ex:recurrence-of-unit-speed-geodesics}
设 $(M,\langle\cdot,\cdot\rangle)$ 为紧 Riemann 流形. 证明对于每个正规球 $B\subset M$ 和每个 $T>0$, 都存在满足 $\|\dot c(t)\|=1$ 的测地线 $c:\mathbb{R}\to M$, 使得 $c(0)\in B$, 并且对某个 $t\geq T$, 有 $c(t)\in B$.
\item\label{ex:hamiltonian-fields-of-canonical-coordinates}
设 $x^1,\ldots,x^n,p_1,\ldots,p_n$ 为 $T^*M$ 上通常的局部坐标. 求 $X_{x^i}$, $X_{p_i}$, $\{x^i,x^j\}$, $\{p_i,p_j\}$ 和 $\{p_i,x^j\}$.
\item\label{ex:leibniz-rule-for-poisson-bracket}
证明对所有 $F,G,H\in C^\infty(T^*M)$, Poisson 括号满足 \textbf{Leibniz 法则}
\[
  \{F,GH\}=\{F,G\}H+\{F,H\}G.
\]
\end{enumerate}
\end{RiemannExercise}
